\documentclass[10pt,a4paper]{amsart}
\usepackage{amsmath,amssymb,amsthm,graphicx}
\usepackage{verbatim}
\newenvironment{reference}{\comment}{\endcomment}
%\newenvironment{reference}{}{}
\newenvironment{slogan}{\comment}{\endcomment}
\newenvironment{history}{\comment}{\endcomment}

% For commutative diagrams we use Xy-pic
\usepackage[all]{xy}

% We use 2cell for 2-commutative diagrams.
\xyoption{2cell}
% Theorem environments.
%
\theoremstyle{plain}
\newtheorem{theorem}[subsection]{Theorem}
\newtheorem{proposition}[subsection]{Proposition}
\newtheorem{lemma}[subsection]{Lemma}

\theoremstyle{definition}
\newtheorem{definition}[subsection]{Definition}
\newtheorem{example}[subsection]{Example}
\newtheorem{exercise}[subsection]{Exercise}
\newtheorem{situation}[subsection]{Situation}

\theoremstyle{remark}
\newtheorem{remark}[subsection]{Remark}
\newtheorem{remarks}[subsection]{Remarks}

\numberwithin{equation}{subsection}

% Macros
%
\def\lim{\mathop{\mathrm{lim}}\nolimits}
\def\colim{\mathop{\mathrm{colim}}\nolimits}
\def\Spec{\mathop{\mathrm{Spec}}}
\def\Hom{\mathop{\mathrm{Hom}}\nolimits}
\def\Ext{\mathop{\mathrm{Ext}}\nolimits}
\def\SheafHom{\mathop{\mathcal{H}\!\mathit{om}}\nolimits}
\def\SheafExt{\mathop{\mathcal{E}\!\mathit{xt}}\nolimits}
\def\Sch{\mathit{Sch}}
\def\Mor{\mathop{\mathrm{Mor}}\nolimits}
\def\Ob{\mathop{\mathrm{Ob}}\nolimits}
\def\Sh{\mathop{\mathit{Sh}}\nolimits}
\def\NL{\mathop{N\!L}\nolimits}
\def\CH{\mathop{\mathrm{CH}}\nolimits}
\def\proetale{{pro\text{-}\acute{e}tale}}
\def\etale{{\acute{e}tale}}
\def\QCoh{\mathit{QCoh}}
\def\Ker{\mathop{\mathrm{Ker}}}
\def\Im{\mathop{\mathrm{Im}}}
\def\Coker{\mathop{\mathrm{Coker}}}
\def\Coim{\mathop{\mathrm{Coim}}}

% Boxtimes
%
\DeclareMathSymbol{\boxtimes}{\mathbin}{AMSa}{"02}

%
% Macros for moduli stacks/spaces
%
\def\QCohstack{\mathcal{QC}\!\mathit{oh}}
\def\Cohstack{\mathcal{C}\!\mathit{oh}}
\def\Spacesstack{\mathcal{S}\!\mathit{paces}}
\def\Quotfunctor{\mathrm{Quot}}
\def\Hilbfunctor{\mathrm{Hilb}}
\def\Curvesstack{\mathcal{C}\!\mathit{urves}}
\def\Polarizedstack{\mathcal{P}\!\mathit{olarized}}
\def\Complexesstack{\mathcal{C}\!\mathit{omplexes}}
% \Pic is the operator that assigns to X its picard group, usage \Pic(X)
% \Picardstack_{X/B} denotes the Picard stack of X over B
% \Picardfunctor_{X/B} denotes the Picard functor of X over B
\def\Pic{\mathop{\mathrm{Pic}}\nolimits}
\def\Picardstack{\mathcal{P}\!\mathit{ic}}
\def\Picardfunctor{\mathrm{Pic}}
\def\Deformationcategory{\mathcal{D}\!\mathit{ef}}
\usepackage{cleveref}
\crefname{originallemma}{original lemma}{original lemmas}
\Crefname{originallemma}{Original lemma}{Original lemmas}
\makeatletter
\newlabel{lemma-isolated-point}{{00PJ}{source}{}{sourcepointer-1}{}}
\newlabel{lemma-isolated-point@cref}{{[originallemma][0][]00PJ}{[1][0][]source}{}{}{}}
\makeatother
\crefname{originallemma}{original lemma}{original lemmas}
\Crefname{originallemma}{Original lemma}{Original lemmas}
\makeatletter
\newlabel{lemma-spec-product}{{00ED}{source}{}{sourcepointer-2}{}}
\newlabel{lemma-spec-product@cref}{{[originallemma][0][]00ED}{[1][0][]source}{}{}{}}
\makeatother
\crefname{originallemma}{original lemma}{original lemmas}
\Crefname{originallemma}{Original lemma}{Original lemmas}
\makeatletter
\newlabel{proposition-dimension-zero-ring}{{00KJ}{source}{}{sourcepointer-3}{}}
\newlabel{proposition-dimension-zero-ring@cref}{{[originallemma][0][]00KJ}{[1][0][]source}{}{}{}}
\makeatother
\title{Quasi-finiteness at a prime is preserved by a surjective change of the four rings}
\author{The Stacks Project Authors}
\date{stacks-project:algebra.tex:local-snapshot-20261012; GNU FDL 1.2 or later}
\begin{document}
\maketitle
\noindent\textit{Editorial note.} All delivered mathematical text is the author's own, copied from the pinned Stacks Project chapter file algebra.tex; no statement or proof step was rewritten, reordered or removed.\par\medskip
\noindent\textit{Editorial note.} Preamble provenance, as required for this source: the frozen snapshot stages the chapter's own preamble file beside it (preamble.tex, 217 lines, 8116 bytes, the same shared Stacks-wide preamble the published chow and etale-cohomology cards use); the first line of the pinned chapter file is \textbackslash{}input\{preamble\}, so no preamble text lives inside the chapter itself. The card takes that file's lines 8-18 (the verbatim package and the three \textbackslash{}comment-based environments reference/slogan/history at lines 9-12, section 13.2, together with the Xy-pic setup) and its lines 154 to the runtime end of the file (the native theorem declarations and mathematical macros). Not one character was changed, no package was added, no macro was defined, and the excerpt serves only as preamble, never as a span coordinate.\par\medskip
\noindent\textit{Editorial note.} The same-chapter pointers are located at their \textbackslash{}label lines inside the pinned chapter file; the cross-chapter pointers are located in the project's own tag list staged in the snapshot (sources/stacks-chow/tags.txt), at the verbatim lines carrying their canonical tags.\par\medskip
\noindent\textit{Editorial note.} The counted statement contains an Xy-pic diagram; the preamble spans of this card run 8-18 of the chapter's own preamble file and therefore include the Xy-pic package and option setup at lines 15-18.\par\medskip
\noindent\textit{Editorial note.} The counted proof refers to same-chapter statements by label, recorded as navigation stubs; no \textbackslash{}cite appears, so the card carries no bibliography.\par\medskip
\noindent\textit{Allowed standard prerequisites (editorial).}\begin{itemize}
\item Isolated points of a fibre give a product decomposition. Invoked at the start of the counted proof (lines 30527-30531). Writes the fibre as a product with a finite factor containing the point.
\item Finiteness and dimension zero detected after a product decomposition. Used at the end of the counted proof (lines 30545-30554). Shows the point stays isolated in the new fibre.
\end{itemize}\medskip
\ifdefined\hypertarget\hypertarget{sourcepointer-1}{}\fi
\noindent\textit{Source reference.} Original lemma 00PJ: sources/stacks-algebra/pinned-original-source.tex, source lines 30323--30323. Navigation pointer for \textbackslash{}ref\{lemma-isolated-point\} used inside the counted span (line 30530); the argument that consumes it is retained in the span.\par\medskip
\ifdefined\hypertarget\hypertarget{sourcepointer-2}{}\fi
\noindent\textit{Source reference.} Original lemma 00ED: sources/stacks-algebra/pinned-original-source.tex, source lines 4039--4039. Navigation pointer for \textbackslash{}ref\{lemma-spec-product\} used inside the counted span (line 30550); the argument that consumes it is retained in the span.\par\medskip
\ifdefined\hypertarget\hypertarget{sourcepointer-3}{}\fi
\noindent\textit{Source reference.} Original lemma 00KJ: sources/stacks-algebra/pinned-original-source.tex, source lines 14632--14632. Navigation pointer for \textbackslash{}ref\{proposition-dimension-zero-ring\} used inside the counted span (line 30553); the argument that consumes it is retained in the span.\par\medskip
\begin{lemma}
\label{lemma-four-rings}
Let
$$
\xymatrix{
S \ar[r] & S' & &
\mathfrak q \ar@{-}[r] & \mathfrak q' \\
R \ar[u] \ar[r] &  R' \ar[u] & &
\mathfrak p \ar@{-}[r] \ar@{-}[u] & \mathfrak p' \ar@{-}[u]
}
$$
be a commutative diagram of rings with primes as indicated.
Assume $R \to S$ of finite type, and $S \otimes_R R' \to S'$ surjective.
If $R \to S$ is quasi-finite at $\mathfrak q$, then
$R' \to S'$ is quasi-finite at $\mathfrak q'$.
\end{lemma}

\begin{proof}
Write $S \otimes_R \kappa(\mathfrak p) = S_1 \times S_2$
with $S_1$ finite over $\kappa(\mathfrak p)$ and such that
$\mathfrak q$ corresponds to a point of $S_1$ as in
Lemma \ref{lemma-isolated-point}. This product decomposition
induces a corresponding product decomposition for any
$S \otimes_R \kappa(\mathfrak p)$-algebra. In particular,
we obtain $S' \otimes_{R'} \kappa(\mathfrak p') = S'_1 \times S'_2$.
Because $S \otimes_R R' \to S'$ is surjective the canonical map
$(S \otimes_R \kappa(\mathfrak p)) \otimes_{\kappa(\mathfrak p)}
\kappa(\mathfrak p') \to S' \otimes_{R'} \kappa(\mathfrak p')$
is surjective and hence
$S_i \otimes_{\kappa(\mathfrak p)} \kappa(\mathfrak p') \to S'_i$
is surjective. It follows that $S'_1$ is finite over $\kappa(\mathfrak p')$.
The map $S' \otimes_{R'} \kappa(\mathfrak p') \to
\kappa(\mathfrak q')$ factors through $S_1'$
(i.e.\ it annihilates the factor $S_2'$)
because the map $S \otimes_R \kappa(\mathfrak p) \to
\kappa(\mathfrak q)$ factors through $S_1$
(i.e.\ it annihilates the factor $S_2$). Thus
$\mathfrak q'$ corresponds to a point of
$\Spec(S_1')$ in the disjoint union decomposition
of the fibre: $\Spec(S' \otimes_{R'} \kappa(\mathfrak p'))
= \Spec(S_1') \amalg \Spec(S_2')$, see
Lemma \ref{lemma-spec-product}.
Since $S_1'$ is finite over a field, it is Artinian ring,
and hence $\Spec(S_1')$ is a finite discrete set.
(See Proposition \ref{proposition-dimension-zero-ring}.)
We conclude $\mathfrak q'$ is isolated in its fibre as
desired.
\end{proof}

\end{document}
